\lesson{6}{от 6 мая 2025 12:45}{Продолжение}

Градиент, который мы только что нашли, показывает \emph{направление скорейшего подъема} (возрастания) функции.

Нам же, если мы хотим найти минимум функции, нужно двигаться вниз...

\begin{definition}[Метод градиентного спуска]
  \emph{Метод градиентного спуска} -- это способ нахождения локального минимума функции в процессе движения в направлении антиградиента.
\end{definition}

\begin{note}
  При этом, если мы сдвинемся на всю величину градиента, то "перескочим"...

  \begin{itemize}
    \item после ста итераций в примере, как веса модели, так и уровень ошибки приблизились к нулю;
      \item более сложные функции потерь потребуют большего количества итераций и более тонкой настройки коэфициента скорости обучения.
  \end{itemize}
\end{note}

\subsection{Промежуточные выводы}

\begin{enumerate}
  \item Градиент показывает скорость изменения многомерной функции по каждому из измерений.
    \item Если "заморозить" все независимые переменные кроме одной и сделать срез, то наклон казательной к кривой среза будет значением частной производной по этому измерению.
      \item Градиент показывает направление скорейшего подъема функции. Двигаясь в обратном направлении, мы придем к минимуму функции.
\end{enumerate}

\section{Глубокое обучение}

\begin{enumerate}
  \item Определения.
    \item DL vs ML
      \item Классы DL.
        \item Струткра модели DL.
          \item Нейросети.
\end{enumerate}

\begin{definition}[Искусственный интелект]
  \emph{Искусственный интелект} -- это способность компьютеров решать задачи при помощи вычислений.
\end{definition}

\emph{Проводить кластерный анализ}. Обученные без учителя нейронные сети лучше всего справляются с задачей кластеризации -- делением большого массива данных на группы. Нейросеть способна выделить критерии отличия и иоделить данные на их основе. А потом с группами уже смогут работать люди, находя нужные тенденции и проводя аналитику. 

\emph{Выявлять аномалии}. Для этого нейросить сначала обучают на нормальных данных, и, когда в реальности она встречается с выбросами, то должна об этом сигнализировать. Так можно быстро обнаруживать необычное поведение покупателей, например резкий рост спроса на конкретный товар.

\emph{Находить ассоциации}. Нейросеть может определить критерии похожести одних объектов на другие и строить связи.

Например, рекомендовать товары в дополнение к купленным.
